% =========================================================
% PAGE 23: SUB-INDICES PERFORMANCE (Figure 3.1) - ARABIC
% =========================================================

% 1. Geometry
% Standard margins
\newgeometry{left=2cm, right=2cm, top=1cm, bottom=2.5cm}

% ---------------------------------------------------------
% SECTION TITLE
% ---------------------------------------------------------
\noindent
{\fontsize{18}{22}\selectfont \textbf{\color{ecestitle} \ECESContent{title_section_iii_ar}{ثالثاً: تقييم الأداء والتوقعات وفقاً للمؤشرات الفرعية}}}
\vspace{1em}
\label{sec:ar_eval_sub}
% ---------------------------------------------------------
% SUBSECTION 3.1
% ---------------------------------------------------------
\noindent
{\fontsize{11}{14}\selectfont \textbf{\color{ecestitle} \ECESContent{title_subindices_31_ar}{3.1 تقييم الأداء}}}

\vspace{0.8em}

% ---------------------------------------------------------
% BODY TEXT
% ---------------------------------------------------------
\noindent
{\fontsize{10}{13}\selectfont
\ECESContent{subindices_perf_p1_ar}{[أدخل: أداء مؤشرات الإنتاج والمبيعات المحلية واستغلال الطاقة الإنتاجية والصادرات والمخزون للشركات الكبيرة خلال الربع — القيم مقارنة بالمستوى المحايد والربعين السابق والمناظر.]}

\vspace{0.8em}

\ECESContent{subindices_perf_p2_ar}{[أدخل: أداء نفس المؤشرات للشركات الصغيرة والمتوسطة خلال الربع — القيم مقارنة بالمستوى المحايد والربعين السابق والمناظر.]}
}

\vspace{0.5em}

% ---------------------------------------------------------
% FIGURE 3.1 (Sub-indices Performance)
% ---------------------------------------------------------
\begin{figure}[h!]
    \centering
    % Contains the Production and Sales Indices chart
    \includegraphics[width=\linewidth]{ch20.png}
\end{figure}

% ---------------------------------------------------------
% SOURCE & NOTES
% ---------------------------------------------------------
\vspace{-0.5em}
\noindent
{\fontsize{8}{10}\selectfont
\ECESContent{subindices_source_pg23_ar}{المصدر: نتائج الاستبيان.\\
* تم عكس مؤشر المخزون السلعي ليعكس الأثر السلبي الناتج عن ارتفاع المخزون السلعي على الشركات. لذلك فكلما كان المؤشر السلعي مرتفعا كان المخزون السلعي منخفض والعكس صحيح.}
}

\clearpage

% =========================================================
% PAGE 24: PRICES, WAGES, INVESTMENT, EMPLOYMENT INDICES - ARABIC
% =========================================================

% 1. Geometry (Persist)

% ---------------------------------------------------------
% SECTION TITLE: CONTINUATION FROM PAGE 23
% ---------------------------------------------------------
\noindent
{\fontsize{10}{13}\selectfont \textbf{\ECESContent{subindices_prices_wages_intro_ar}{[أدخل: العنوان الرئيسي البارز — ملخص أداء مؤشري أسعار المنتجات النهائية والأجور مقارنة بالمستوى المحايد.]}}}

\vspace{0.5em}

\noindent
{\fontsize{10}{13}\selectfont
\ECESContent{subindices_prices_wages_p1_ar}{[أدخل: تحليل مؤشر أسعار المنتجات النهائية ومؤشر الأجور لكافة الشركات خلال الربع — القيم مقارنة بالمستوى المحايد والربع السابق (الشكل 3-2).]}
}

\vspace{1em}

% ---------------------------------------------------------
% FIGURE 3.2 (Prices/Wages)
% ---------------------------------------------------------
\begin{figure}[h!]
    \centering
    % Price and Production Cost Indices
    \includegraphics[width=\linewidth]{ch21.png}
\end{figure}

\vspace{-0.5em}
\noindent
{\fontsize{8}{10}\selectfont
\ECESContent{subindices_source_pg24_1_ar}{المصدر: نتائج الاستبيان.\\
* تم عكس مؤشر أسعار المدخلات ليعكس الأثر السلبي الناتج عن ارتفاع الأسعار على المؤشر الإجمالي. لذلك فكلما كان المؤشر منخفضا كان الارتفاع السعري أكبر والعكس صحيح.}
}

\vspace{1.5em}

% ---------------------------------------------------------
% SECTION 2: INVESTMENT AND EMPLOYMENT
% ---------------------------------------------------------
\noindent
{\fontsize{11}{13}\selectfont \textbf{\ECESContent{subindices_inv_emp_intro_ar}{[أدخل: العنوان الرئيسي البارز — ملخص أداء مؤشري الاستثمار والتوظيف مقارنة بالمستوى المحايد.]}}}

\vspace{0.6em}

\noindent
{\fontsize{10}{13}\selectfont
\ECESContent{subindices_inv_emp_large_ar}{[أدخل: أداء مؤشري الاستثمار والتوظيف للشركات الكبيرة خلال الربع — القيم مقارنة بالمستوى المحايد والربعين السابق والمناظر.]}

\vspace{0.6em}

\ECESContent{subindices_inv_emp_sme_ar}{[أدخل: أداء مؤشري الاستثمار والتوظيف للشركات الصغيرة والمتوسطة خلال الربع — القيم مقارنة بالمستوى المحايد والربعين السابق والمناظر (الشكل 3-3).]}
}

\vspace{0.5em}

% ---------------------------------------------------------
% FIGURE 3.3 (Invest/Employ)
% ---------------------------------------------------------
\begin{figure}[h!]
    \centering
    % Investment and Employment Indices
    \includegraphics[width=\linewidth]{ch22.png}
\end{figure}

\vspace{-0.5em}
\noindent
{\fontsize{8}{10}\selectfont
\ECESContent{subindices_source_pg24_2_ar}{المصدر: نتائج الاستبيان.}
}

\clearpage

% =========================================================
% PAGE 25: PERFORMANCE EXPECTATIONS (Figure 3.4) - ARABIC
% =========================================================

% 1. Geometry (Persist)


% ---------------------------------------------------------
% SECTION HEADER
% ---------------------------------------------------------
\noindent
{\fontsize{11}{14}\selectfont \textbf{\color{ecestitle} \ECESContent{title_subindices_32_ar}{3.2 توقعات الأداء}}}

\vspace{0.8em}

% ---------------------------------------------------------
% INTRO TEXT
% ---------------------------------------------------------
\noindent
{\fontsize{10}{13}\selectfont
\textbf{\ECESContent{subindices_exp_intro_ar}{[أدخل: العنوان الرئيسي البارز — ملخص توقعات الأداء لكافة الشركات خلال الربع القادم مقارنة بالمستوى المحايد.]}}}

\vspace{0.8em}

\ECESContent{subindices_exp_large_ar}{[أدخل: توقعات الشركات الكبيرة لمؤشري الإنتاج واستغلال الطاقة والمبيعات المحلية والصادرات والمخزون خلال الربع (\NextQuarterText) مقارنة بالربعين السابق والمناظر.]}

\vspace{0.8em}

\ECESContent{subindices_exp_sme_ar}{[أدخل: توقعات الشركات الصغيرة والمتوسطة لنفس المؤشرات خلال الربع (\NextQuarterText) مقارنة بالربعين السابق والمناظر (الشكل 3-4).]}
}

\vspace{0.5em}

% ---------------------------------------------------------
% FIGURE 3.4 (Expectations Size)
% ---------------------------------------------------------
\begin{figure}[h!]
    \centering
    % Production and Sales Indices by Firm Size - Outlook
    \includegraphics[width=\linewidth]{ch23.png}
\end{figure}

% ---------------------------------------------------------
% SOURCE & NOTES
% ---------------------------------------------------------
\vspace{-0.5em}
\noindent
{\fontsize{8}{10}\selectfont
\ECESContent{subindices_source_pg25_ar}{المصدر: نتائج الاستبيان.\\
* تم عكس مؤشر المخزون السلعي ليعكس الأثر السلبي الناتج عن ارتفاع المخزون السلعي على الشركات. لذلك فكلما كان المؤشر منخفضا كان المخزون السلعي مرتفعا والعكس صحيح.}
}

\vspace{1.5em}

% ---------------------------------------------------------
% BOTTOM TEXT
% ---------------------------------------------------------
\noindent
{\fontsize{10}{13}\selectfont \textbf{\ECESContent{subindices_exp_prices_wages_intro_ar}{[أدخل: العنوان الرئيسي البارز — توقعات مؤشري أسعار المنتجات النهائية والأجور مقارنة بالمستوى المحايد.]}}}

\vspace{0.8em}

\noindent
{\fontsize{10}{13}\selectfont
\ECESContent{subindices_exp_prices_p1_ar}{[أدخل: تحليل توقعات مؤشر أسعار المنتجات النهائية للشركات الكبيرة — مقارنة بالربع السابق والمناظر ومؤشر أسعار المدخلات.]}

\vspace{0.8em}

\ECESContent{subindices_exp_wages_sme_ar}{[أدخل: توقعات مؤشر الأجور للشركات الصغيرة والمتوسطة والكبيرة للربع القادم — مقارنة بالمستوى المحايد والربعين السابق والمناظر (الشكل 3-5).]}
}

\clearpage

% =========================================================
% PAGE 26: PRICES, WAGES, INVESTMENT OUTLOOK - ARABIC
% =========================================================


% ---------------------------------------------------------
% FIGURE 3.5 (Expectations Prices)
% ---------------------------------------------------------
\begin{figure}[h!]
    \centering
    % Prices and Production Cost Indices by Firm Size - Outlook
    \includegraphics[width=\linewidth]{ch24.png}
\end{figure}

% ---------------------------------------------------------
% SOURCE & NOTES FOR FIG 3.5
% ---------------------------------------------------------
\vspace{-0.5em}
\noindent
{\fontsize{8}{10}\selectfont
\ECESContent{subindices_source_pg26_1_ar}{المصدر: نتائج الاستبيان.\\
* تم عكس مؤشر أسعار المدخلات ليعكس الأثر السلبي الناتج عن ارتفاع الأسعار على المؤشر الإجمالي. لذلك فكلما كان المؤشر منخفضا كان الارتفاع السعري أكبر والعكس صحيح.}
}

\vspace{2em}

% ---------------------------------------------------------
% MIDDLE SECTION: INVESTMENT AND EMPLOYMENT
% ---------------------------------------------------------
\noindent
{\fontsize{10}{13}\selectfont \textbf{\ECESContent{subindices_exp_inv_emp_intro_ar}{[أدخل: العنوان الرئيسي البارز — توقعات مؤشري الاستثمار والتوظيف مقارنة بالمستوى المحايد.]}}}

\vspace{0.8em}

\noindent
{\fontsize{10}{13}\selectfont
\ECESContent{subindices_exp_inv_emp_p1_ar}{[أدخل: تحليل توقعات مؤشري الاستثمار والتوظيف لكافة الشركات خلال الربع القادم — مقارنة بالمستوى المحايد والربعين السابق والمناظر (الشكل 3-6).]}
}

\vspace{0.5em}

% ---------------------------------------------------------
% FIGURE 3.6 (Expectations Invest/Employ)
% ---------------------------------------------------------
\begin{figure}[h!]
    \centering
    % Investment and Employment Indices by Firm Size - Outlook
    \includegraphics[width=\linewidth]{ch25.png}
\end{figure}

% ---------------------------------------------------------
% SOURCE FOR FIG 3.6
% ---------------------------------------------------------
\vspace{-0.5em}
\noindent
{\fontsize{8}{10}\selectfont
\ECESContent{subindices_source_pg26_2_ar}{المصدر: نتائج الاستبيان.}
}

\clearpage
